\documentclass[11pt,a4paper]{article}
\usepackage[utf8]{inputenc}
\usepackage[spanish]{babel}
\usepackage{amsmath,amssymb}
\usepackage{booktabs}
\usepackage{geometry}
\geometry{margin=2.5cm}
\usepackage{xcolor}
\definecolor{navyblue}{RGB}{15,34,64}
\usepackage{hyperref}
\hypersetup{colorlinks=true,linkcolor=navyblue,urlcolor=navyblue}

\title{\textbf{Guía Teórica y Econométrica: Modelos OLS, HC1 y Newey-West (HAC)}\\\large\textmd{Corrección de Errores Estándar en Regresiones CAPM}}
\author{\textbf{Julián Duarte} -- DEFI 2}
\date{Septiembre de 2026}

\begin{document}
\maketitle

\section{El Problema de los Errores Estándar en Finanzas}
Al estimar la regresión CAPM:
\begin{equation}
R_{i,t} - R_{f,t} = \alpha_i + \beta_i (R_{m,t} - R_{f,t}) + \varepsilon_{i,t}
\end{equation}
el método MCO (OLS) supone que los errores $\varepsilon_{i,t}$ tienen varianza constante (homocedasticidad) y son independientes a lo largo del tiempo.

En la realidad de los mercados financieros:
\begin{enumerate}
    \item \textbf{Heterocedasticidad:} Hay periodos de calma alternados con rachas de turbulencia extrema (\textit{volatility clustering}).
    \item \textbf{Autocorrelación:} La información y el volumen se incorporan progresivamente a lo largo de varios días.
\end{enumerate}

\section{Comparativa de Estimadores}
\subsection{1. OLS Estándar}
Asume varianza homocedástica: $\widehat{\text{Var}}_{\text{OLS}}(\hat{\beta}) = s^2 (X^T X)^{-1}$.

\subsection{2. HC1 (White con Grados de Libertad)}
Permite varianzas distintas en cada residuo $e_i$:
\begin{equation}
\widehat{\text{Var}}_{\text{HC1}}(\hat{\beta}) = \frac{n}{n-k} (X^T X)^{-1} \left( \sum_{i=1}^n e_i^2 x_i x_i^T \right) (X^T X)^{-1}
\end{equation}

\subsection{3. Newey-West (HAC)}
Corrige heterocedasticidad Y autocorrelación mediante el núcleo de Bartlett ($L=5$):
\begin{equation}
\widehat{\text{Var}}_{\text{HAC}}(\hat{\beta}) = (X^T X)^{-1} \hat{\Omega}_{\text{HAC}} (X^T X)^{-1}
\end{equation}
\begin{equation}
\hat{\Omega}_{\text{HAC}} = \hat{\Gamma}_0 + \sum_{j=1}^{L} \left(1 - \frac{j}{L+1}\right) (\hat{\Gamma}_j + \hat{\Gamma}_j^T)
\end{equation}

\end{document}
